\begin{textsample}{NEG Dim 4 – persona\_gemini – Score -1.00 – t380\_gemini.txt}  \label{ex:f4_neg_007}
Greenpeace International, in coalition with 350 other organizations, has denounced a recent agreement between the Food and Agriculture Organization of the United Nations ( FAO ) and CropLife International.
CropLife International serves as a lobby for major pesticide and genetically engineered \textbf{seed} corporations, including Bayer, BASF, and Syngenta.
The formal goal of this partnership is to transform agri-food systems, promote rural development, and boost farmer resilience to climate change through investment and \textbf{innovation}.
However, critics argue that this agreement empowers profit-driven multinational corporations, granting them undue legitimacy and control over global food systems.
This move ignores vital approaches such as agroecology, food sovereignty, and the recommendations laid out by the 2009 International Assessment of Agricultural Knowledge, Science and Technology for Development ( IAASTD ).
For the sake of a healthy planet, the FAO must prioritize ecological farming over corporate interests.

% matched lemmas: innovation, seed
\end{textsample}
